\documentclass[11pt]{article}
\usepackage[a4paper,margin=25mm]{geometry}
\usepackage{amsmath,amssymb,amsthm}
\usepackage[hidelinks]{hyperref}
\newtheorem{theorem}{Theorem}
\title{Diagonal solutions do not give the proposed VMVT coefficient\\
\large Disproof of Conjecture 00000004157}
\author{gaochengzhecpu}
\date{}
\begin{document}
\maketitle
\begin{abstract}
For $k=1,s=2$, inside the stated range $s\leq k(k+1)$, the exact Vinogradov mean value is $(2N^3+N)/3$, whereas the diagonal count is $2N^2-N$. Thus the true cubic coefficient is $2/3$, and the diagonal contribution on that same scale is zero. The diagonal count's own leading coefficient is $2!=2$, which also differs from $2/3$. As a supplementary degree-two obstruction, an elementary counting argument gives $J_{6,2}(N)\geq N^9/49$, while there are at most $6!N^6$ diagonal solutions. The Lean project proves these facts for the actual finite solution sets.
\end{abstract}

\section*{Statement and conventions}
The source defines the optimal constant as the main-term coefficient of $J_{s,k}(X)$. It asserts that this coefficient always equals the number of diagonal solutions, expressed by factorials, with validity exactly for $s\leq k(k+1)$. It imposes no restriction excluding $k=1$.

For an integer $N\geq1$, use the standard unweighted count
\[
 J_{s,k}(N)=\#\left\{(\boldsymbol x,\boldsymbol y)\in\{1,\ldots,N\}^{2s}:
 \sum_{i=1}^s x_i^j=\sum_{i=1}^s y_i^j\quad(1\leq j\leq k)\right\}.
\]
By a diagonal solution we mean an ordered pair for which $\boldsymbol y$ is a permutation of $\boldsymbol x$; each pair is counted once, even when repeated coordinates allow several permutations. Write its count as $D_s(N)$. These are the usual solution counts for the Vinogradov mean value. Disagreement along the integer values of $X$ suffices to refute an asserted general asymptotic.

\section*{An exact counterexample to the coefficient claim}
\begin{theorem}
For every integer $N\geq1$,
\[
 J_{2,1}(N)=\frac{2N^3+N}{3},\qquad D_2(N)=2N^2-N.
\]
Consequently
\[
 \lim_{N\to\infty}\frac{J_{2,1}(N)}{N^3}=\frac23,
 \qquad \lim_{N\to\infty}\frac{D_2(N)}{N^3}=0,
 \qquad \lim_{N\to\infty}\frac{D_2(N)}{N^2}=2.
\]
\end{theorem}
\begin{proof}
The only equation is $a+b=c+d$. Translating all four coordinates by $-1$ changes the domain to $\{0,\ldots,N-1\}$ without changing this equation. Classify a solution by $\delta=a-c=d-b$.

When $\delta=0$, there are $N^2$ solutions. If $1\leq\delta\leq N-1$, choose $c,b\in\{0,\ldots,N-1-\delta\}$ freely; the remaining coordinates are uniquely $a=c+\delta$ and $d=b+\delta$. Thus there are $(N-\delta)^2$ such solutions. Interchanging the left and right pairs gives the same count for $-\delta$. Hence
\[
 J_{2,1}(N)=N^2+2\sum_{m=0}^{N-1}m^2
 =N^2+\frac{N(N-1)(2N-1)}3
 =\frac{2N^3+N}{3}.
\]
For the diagonal count the two possible right pairs are $(a,b)$ and $(b,a)$. Each choice gives $N^2$ ordered solutions, and the two families intersect exactly when $a=b$, in $N$ solutions. Therefore $D_2(N)=2N^2-N$. Dividing these exact formulas by the displayed powers of $N$ proves the limits.
\end{proof}

Here $s=2=k(k+1)$ for $k=1$, so the example belongs to the claimed range. On the actual main-term scale $N^3$, the coefficient of $J_{2,1}$ is $2/3$ while the diagonal contribution is zero. If the wording is instead read as asserting equality with the factorial coefficient of the diagonal count itself, that coefficient is $2!=2$, again unequal to $2/3$. Thus the counterexample directly concerns the coefficient, not only the existence of some non-diagonal solutions.

\section*{A supplementary obstruction with degree two}
There is also a failure of diagonal main-term dominance at $k=2,s=6$, again with $s=k(k+1)$. Map each $\boldsymbol x\in\{1,\ldots,N\}^6$ to its genuine moment pair
\[
 F(\boldsymbol x)=\left(\sum_{i=1}^6 x_i,\ \sum_{i=1}^6 x_i^2\right).
\]
Both entries are nonnegative integers, at most $6N$ and $6N^2$, respectively. Thus they lie in a box with
\[
 B_N=(6N+1)(6N^2+1)\leq49N^3
\]
elements. If $r_b$ is the number of tuples in the fiber over $b$, then
\[
 \sum_b r_b=N^6,\qquad J_{6,2}(N)=\sum_b r_b^2.
\]
Cauchy--Schwarz gives
\[
 N^{12}\leq B_N J_{6,2}(N)\leq49N^3J_{6,2}(N),
 \qquad J_{6,2}(N)\geq\frac{N^9}{49}.
\]
Choosing the left tuple and one of the $6!$ permutations covers every diagonal pair, possibly with repetitions. Hence $D_6(N)\leq720N^6$. It follows that $J_{6,2}(N)$ is not $O(N^6)$, and in particular cannot be asymptotic to $D_6(N)$. This is a supplementary refutation of diagonal dominance within the claimed range; it does not purport to compute the exact degree-two asymptotic coefficient. The preceding degree-one example already supplies an exact coefficient counterexample to the source as written.

\section*{Formal correspondence and verification}
The Lean project works with actual finite types and their cardinalities. For $k=1,s=2$, \texttt{LinearSolutions} is the subtype of four entries in \texttt{Fin N} satisfying $a+b=c+d$. The theorem \texttt{linear\_equation\_iff\_standard} proves that adding one to all entries gives the standard equation on $\{1,\ldots,N\}$. The map \texttt{encodeLinear} implements the three difference cases in the proof; its injectivity and surjectivity are proved before its cardinality is used. The diagonal set is the union of the identity and transposition families, with its intersection counted exactly.

The theorems \texttt{J21\_exact} and \texttt{D21\_exact} establish the two polynomial formulas. The actual real limits, including the zero diagonal contribution on the cubic scale, are proved using \texttt{Filter.Tendsto}. They are combined with the range check in \texttt{explicit\_\allowbreak coefficient\_\allowbreak counterexample}.

For $k=2,s=6$, \texttt{momentMap} sends actual six-tuples to the two integer power sums. The theorem \texttt{momentMap\_eq\_iff} identifies collisions with the two Vinogradov equations. An explicit equivalence with a disjoint union of pairs of fibers proves the collision formula; finite Cauchy--Schwarz gives the lower bound. The diagonal type consists of actual permuted tuple pairs, and these are proved to satisfy both equations. The project then proves the $720N^6$ bound and negates asymptotic equivalence with the full solution count.

The project is pinned to Lean 4.19.0 and a fixed Mathlib commit. No custom axioms, proof placeholders, or numerical oracles are used. Printed theorem dependencies are limited to the standard \texttt{propext}, \texttt{Classical.choice}, and \texttt{Quot.sound}. No external theorem about the VMVT asymptotic is assumed.

\section*{Source}
\href{https://github.com/The-Last-Math-Competition/The-Last-Math-Competition/blob/main/conjectures/00000004157.md}{The Last Math Competition, Conjecture 00000004157}. The submission preserves the exact bilingual source and records its commit and SHA-256 separately.
\end{document}
